% ============================================================
% secciones/02_marco_teorico.tex
% ============================================================

\section{Marco Teórico}
\label{sec:marco}

\subsection{Modelo compuesto frecuencia-severidad}

Sea $S$ la pérdida agregada de una póliza durante un periodo de vigencia. El modelo actuarial estándar descompone $S$ como:

\begin{equation}
  S = \sum_{i=1}^{N} X_i
  \label{eq:modelo_compuesto}
\end{equation}

donde $N$ denota el número de reclamaciones (frecuencia) y $X_i$ el monto de la $i$-ésima reclamación (severidad), asumiendo que $N$ y los $X_i$ son independientes.

\begin{definition}[Prima pura]
La \textbf{prima pura} o \textbf{prima de riesgo} se define como la esperanza de la pérdida agregada:
\begin{equation}
  \Pi = \mathbb{E}[S] = \mathbb{E}[N] \cdot \mathbb{E}[X]
  \label{eq:prima_pura}
\end{equation}
\end{definition}

\subsection{Modelos de frecuencia: distribución Poisson}

Para la componente de frecuencia, se asume $N \sim \text{Poisson}(\lambda)$ con función de masa de probabilidad:

\begin{equation}
  P(N = k) = \frac{e^{-\lambda} \lambda^k}{k!}, \quad k = 0, 1, 2, \ldots
\end{equation}

En el marco GLM con liga logarítmica, la tasa de siniestralidad se modela como:

\begin{equation}
  \log(\lambda_i) = \log(e_i) + \boldsymbol{x}_i^\top \boldsymbol{\beta}
\end{equation}

donde $e_i$ es la exposición (tiempo en riesgo) de la $i$-ésima póliza y $\boldsymbol{x}_i$ el vector de covariables \citep{McCullagh1989}.

\subsection{Modelos de severidad: distribución Gamma}

Para la severidad, dado que $X > 0$, se utiliza la distribución Gamma con función de densidad:

\begin{equation}
  f(x; \alpha, \beta) = \frac{\beta^\alpha}{\Gamma(\alpha)} x^{\alpha-1} e^{-\beta x}, \quad x > 0
\end{equation}

Con liga logarítmica: $\log(\mu_i) = \boldsymbol{x}_i^\top \boldsymbol{\gamma}$, donde $\mu_i = \alpha_i / \beta_i$.

\subsection{Gradient Boosting para datos actuariales}

El algoritmo XGBoost \citep{Chen2016} construye un ensamble de árboles de decisión de forma aditiva:

\begin{equation}
  \hat{y}_i^{(t)} = \hat{y}_i^{(t-1)} + f_t(\boldsymbol{x}_i)
\end{equation}

minimizando en cada iteración $t$ el objetivo regularizado:

\begin{equation}
  \mathcal{L}^{(t)} = \sum_{i=1}^{n} \ell\!\left(y_i,\, \hat{y}_i^{(t)}\right) + \Omega(f_t)
\end{equation}

donde $\Omega(f_t) = \gamma T + \frac{1}{2}\lambda \|\boldsymbol{w}\|^2$ penaliza la complejidad del árbol.

Para la frecuencia se utiliza la función de pérdida de Poisson ($\ell = -\log P(N=k|\lambda)$) y para la severidad la pérdida de Gamma.
